\chapter{Marco Teórico}\label{cap2:marcos}

\section{Conceptos y teorías fundamentales}

Esta sección establece los fundamentos teóricos sobre los cuales se construye
el modelo PINN-SIREN para la simulación de \textit{speckle} óptico.

\subsection{Ecuación de onda de Helmholtz}

La propagación de luz coherente monocromática en un medio homogéneo e isótropo
está gobernada por la ecuación de onda de Helmholtz escalar
\parencite{born1999principles, jin2014finite, goodman2007speckle}:
%
\begin{equation}
    \nabla^2 E(\mathbf{r}) + k^2 E(\mathbf{r}) = 0,
    \label{eq:helmholtz}
\end{equation}
%
donde $E(\mathbf{r}) \in \mathbb{C}$ es la amplitud compleja del campo eléctrico,
$k = 2\pi/\lambda$ es el número de onda y $\lambda$ la longitud de onda del
láser. En dos dimensiones espaciales ($\mathbf{r} = (x,y)$):
%
\begin{equation}
    \frac{\partial^2 E}{\partial x^2} + \frac{\partial^2 E}{\partial y^2}
    + k^2 E = 0.
    \label{eq:helmholtz2d}
\end{equation}

Dado que las redes neuronales operan sobre números reales, el campo complejo se
descompone en sus partes real e imaginaria: $E = E_\text{real} + iE_\text{imag}$.
Ambas componentes satisfacen la ecuación~\eqref{eq:helmholtz2d} de forma
independiente:
%
\begin{equation}
    \nabla^2 E_\text{real} + k^2 E_\text{real} = 0, \qquad
    \nabla^2 E_\text{imag} + k^2 E_\text{imag} = 0.
    \label{eq:helmholtz_split}
\end{equation}

La intensidad física observable del \textit{speckle} se calcula como:
%
\begin{equation}
    I(x,y) = |E|^2 = E_\text{real}^2 + E_\text{imag}^2.
    \label{eq:intensidad}
\end{equation}

\subsection{Speckle óptico completamente desarrollado}

Cuando una superficie rugosa con variación de fase $\Delta\psi \gg 2\pi$ es
iluminada por luz coherente, las contribuciones de los distintos puntos de la
superficie interfieren de forma aleatoria. Por el Teorema del Límite Central,
el campo complejo resultante sigue una distribución gaussiana circular compleja
\parencite{goodman2007speckle}:
%
\begin{equation}
    E_\text{real},\; E_\text{imag} \sim \mathcal{N}(0,\, \sigma_E^2),
    \quad \text{independientes.}
    \label{eq:campo_gaussiano}
\end{equation}

La intensidad $I = |E|^2$ sigue entonces una distribución exponencial negativa:
%
\begin{equation}
    p(I) = \frac{1}{\langle I \rangle}
    \exp\!\left(-\frac{I}{\langle I \rangle}\right), \quad I \geq 0,
    \label{eq:pdf_intensidad}
\end{equation}
%
con valor esperado $\langle I \rangle = 2\sigma_E^2$ y desviación estándar
$\sigma_I = \langle I \rangle$. El \textit{contraste del speckle} se define como:
%
\begin{equation}
    C = \frac{\sigma_I}{\langle I \rangle} = 1,
    \label{eq:contraste}
\end{equation}
%
criterio que vale exactamente 1 para \textit{speckle} completamente desarrollado
y disminuye cuando el campo pierde coherencia o la rugosidad es insuficiente
\parencite{goodman2007speckle}.

\subsection{Redes Neuronales Informadas por Física (PINNs)}

El enfoque de resolver EDPs con redes neuronales fue pionero por
\textcite{lagaris1998ann}, quienes demostraron que redes feedforward podían
aproximar soluciones de ecuaciones diferenciales ordinarias y parciales
mediante una función de prueba que satisface las condiciones de frontera
exactamente. Este trabajo establece el precedente directo de las PINNs modernas.
La capacidad de aproximación universal de las redes feedforward fue formalizada
por \textcite{hornik1989universal}.

Las PINNs modernas, introducidas por \textcite{raissi2019physics} y revisadas
extensamente en \textcite{karniadakis2021physics} y \textcite{cuomo2022scientific},
aproximan la solución de una EDP mediante una red neuronal
$\mathcal{N}_\theta(\mathbf{r})$ cuyos parámetros $\theta$ se optimizan
minimizando una función de pérdida compuesta:
%
\begin{equation}
    \mathcal{L}(\theta) = \mathcal{L}_\text{datos}(\theta)
    + \lambda_\text{fís}\,\mathcal{L}_\text{física}(\theta),
    \label{eq:loss_total}
\end{equation}
%
donde $\mathcal{L}_\text{datos}$ penaliza el error en las condiciones de
frontera y $\mathcal{L}_\text{física}$ penaliza la violación de la EDP en el
interior del dominio. Para la ecuación de Helmholtz 2D con campo complejo:
%
\begin{align}
    \mathcal{L}_\text{datos}(\theta) &=
    \frac{1}{N_b}\sum_{j=1}^{N_b}
    \left\|\mathcal{N}_\theta(\mathbf{r}_j^b) - E(\mathbf{r}_j^b)\right\|^2,
    \label{eq:loss_datos} \\[4pt]
    \mathcal{L}_\text{física}(\theta) &=
    \frac{1}{N_c}\sum_{i=1}^{N_c}
    \bigl[R_\text{real}(\mathbf{r}_i^c)^2 + R_\text{imag}(\mathbf{r}_i^c)^2\bigr],
    \label{eq:loss_fisica}
\end{align}
%
donde $R_\text{real} = \nabla^2 E_\text{real} + k^2 E_\text{real}$ y
$R_\text{imag} = \nabla^2 E_\text{imag} + k^2 E_\text{imag}$ son los residuos
de Helmholtz, $\{\mathbf{r}_j^b\}$ son los puntos de frontera y
$\{\mathbf{r}_i^c\}$ los puntos de colocación interiores. Las derivadas
$\nabla^2 E$ se obtienen mediante diferenciación automática
\parencite{baydin2018autodiff}.

\subsection{Redes SIREN y sesgo espectral}

Las activaciones estándar (ReLU, tanh) presentan un sesgo espectral
\parencite{rahaman2019spectral}: aprenden frecuencias bajas antes que las altas.
Para la ecuación de Helmholtz con $k = 2\pi$, esto produce derivadas de
segundo orden inestables y convergencia lenta. Adicionalmente, ReLU tiene
segunda derivada nula, lo que impide computar el Laplaciano mediante
diferenciación automática.

\textcite{sitzmann2020siren} proponen las redes SIREN
(\textit{Sinusoidal Representation Networks}), que utilizan la activación
$\phi(z) = \sin(\omega_0 z)$ con una inicialización específica:
%
\begin{equation}
    \mathbf{h}_\ell = \sin\!\bigl(\omega_0\,
    (\mathbf{W}_\ell\,\mathbf{h}_{\ell-1} + \mathbf{b}_\ell)\bigr),
    \quad \ell = 1, \ldots, L,
    \label{eq:siren}
\end{equation}
%
donde las matrices $\mathbf{W}_\ell$ se inicializan según
$\mathcal{U}(-\sqrt{6/d},\, \sqrt{6/d})$ para capas intermedias y
$\mathcal{U}(-1/n_\text{in},\, 1/n_\text{in})$ para la primera capa. Esta
inicialización garantiza que las distribuciones de activación sean
aproximadamente uniformes en $[-1,1]$, condición necesaria para la estabilidad
de derivadas de orden arbitrario. La derivada de cualquier orden de SIREN es
también una red SIREN, lo que la convierte en la elección natural para resolver
EDPs con operadores diferenciales de orden superior.

Una alternativa relacionada para mitigar el sesgo espectral es el mapeo de
características de Fourier (\textit{Fourier feature mapping}) propuesto por
\textcite{tancik2020fourier}, que proyecta las coordenadas de entrada a través
de funciones seno/coseno como una capa de preprocesamiento antes de una MLP
estándar. SIREN incorpora este mismo principio de codificación periódica
directamente en la función de activación de cada capa, en lugar de aplicarlo
como una transformación previa de las coordenadas de entrada, lo que le
permite propagar la periodicidad a través de toda la red y facilita el cálculo
de derivadas de orden superior mediante diferenciación automática.

\subsection{Diferenciación automática}

La diferenciación automática (\textit{autodiff}) permite calcular derivadas de
cualquier orden con precisión de máquina mediante la aplicación sucesiva de la
regla de la cadena sobre el grafo de operaciones del programa
\parencite{baydin2018autodiff}. En el contexto de las PINNs, \textit{autodiff}
actúa como el motor que permite evaluar el Laplaciano $\nabla^2$ directamente
sobre las salidas de la red neuronal respecto a sus coordenadas de entrada, sin
los errores de truncamiento de las diferencias finitas.

\subsection{Muestreo Latino Hipercúbico (LHS)}

El Muestreo Latino Hipercúbico \parencite{mckay1979lhs} es una estrategia de
muestreo estratificado diseñada para generar una distribución de puntos de
colocación que cubra de manera cuasi-óptima el dominio computacional $\Omega$.
A diferencia del muestreo aleatorio simple, el LHS divide cada dimensión en
$N_c$ estratos iguales y coloca exactamente un punto por estrato, garantizando
cobertura uniforme sin los patrones redundantes de una malla cartesiana. Esto
permite a la red capturar la física de la ecuación de Helmholtz con una densidad
de puntos significativamente menor, optimizando el uso de memoria GPU y
reduciendo el tiempo de convergencia \parencite{cuomo2022scientific}.

\subsection{Optimización híbrida Adam + L-BFGS}

El entrenamiento sigue una estrategia bifásica de Adam seguido de L-BFGS,
práctica extendida en la literatura de PINNs \parencite{cuomo2022scientific}:

\begin{enumerate}
    \item \textbf{Fase 1 (Adam):} Optimizador estocástico de primer orden
    \parencite{kingma2015adam} con tasa de aprendizaje $\eta = 10^{-3}$,
    paciencia de 800 épocas sin mejora mayor que $\delta = 10^{-7}$ y un
    máximo de 15{,}000 épocas. Adam realiza exploración global del espacio
    de parámetros y proporciona un punto de inicio próximo al mínimo global.

    \item \textbf{Fase 2 (L-BFGS):} Optimizador cuasi-Newton de segundo orden
    (\textit{Limited-memory Broyden-Fletcher-Goldfarb-Shanno})
    \parencite{nocedal2006numerical} con búsqueda de línea \textit{strong Wolfe},
    máximo 1{,}000 iteraciones y tamaño de historial 100. L-BFGS aprovecha la
    información de curvatura para realizar un ajuste fino de alta precisión
    desde el punto entregado por Adam.
\end{enumerate}

Esta combinación supera a Adam puro en precisión final y a L-BFGS puro en
robustez ante mínimos locales.

% ──────────────────────────────────────────────────────────────────────────────
\section{Literatura relacionada}

El uso de redes neuronales para resolver EDPs fue revitalizado por
\textcite{raissi2019physics}, quienes demostraron la capacidad de las PINNs
para resolver una amplia gama de problemas físicos sin necesidad de mallado.
El análisis teórico de por qué las PINNs pueden fallar en el entrenamiento fue
abordado por \textcite{wang2022ntk} mediante la teoría del núcleo tangente
neuronal (NTK), estableciendo las condiciones bajo las cuales el entrenamiento
converge.

En el campo de las ecuaciones de onda, \textcite{moseley2020wave} aplicaron
PINNs para resolver la ecuación de onda acústica 2D en modelos de velocidad
variables, demostrando la viabilidad del enfoque en dominios inhomogéneos.
Para la ecuación de Helmholtz específicamente, \textcite{alkhalifah2021wavefield}
propusieron un marco de aprendizaje automático para obtener soluciones del campo
de onda en sísmica, usando la ecuación de Helmholtz como restricción en la
función de pérdida.

En el campo de la óptica y el electromagnetismo, \textcite{fang2020physics}
aplicaron PINNs para resolver problemas de dispersión de ondas en medios
inhomogéneos 2D, mientras que \textcite{chen2020physics} las utilizaron para
el diseño inverso en nanofotónica. \textcite{lu2021deepxde} desarrollaron la
biblioteca DeepXDE que sistematiza la implementación de PINNs para diversas EDPs.

Un antecedente relevante es el de \textcite{schoder2024helmholtz}, quienes
realizaron un estudio de factibilidad resolviendo la ecuación de Helmholtz
mediante PINNs para propagación de ondas acústicas, demostrando la viabilidad
de la técnica en entornos complejos: reportaron un error relativo $L^2$ de
$0.0249$ (aproximadamente $2.5\%$) entre su solución PINN y una malla FEM de
referencia fina. Este resultado es indicativo del orden de magnitud de
precisión alcanzable con PINNs, aunque no constituye una línea base de
comparación directa para la presente investigación: su experimento es en 3D
(no 1D/2D), su dominio físico es acústico (no óptico), y su métrica de
referencia es una malla FEM (no una solución analítica exacta).

La brecha en el conocimiento identificada es la siguiente: ningún trabajo
previo ha aplicado sistemáticamente la arquitectura SIREN \parencite{sitzmann2020siren}
(con su ventaja intrínseca para representar funciones de alta frecuencia)
a la simulación del \textit{speckle} óptico como fenómeno estadístico complejo.
Esta investigación cierra dicha brecha.

% ──────────────────────────────────────────────────────────────────────────────
\section{Marco tecnológico}

\subsection{Infraestructura de hardware}

El entrenamiento de los modelos PINN-SIREN se realizó sobre la siguiente
infraestructura:

\begin{itemize}
    \item \textbf{GPU:} NVIDIA GeForce RTX 5050 (8.5 GB VRAM), con núcleos
    CUDA para acelerar la diferenciación automática de segundo orden requerida
    por el cálculo del Laplaciano.
    \item \textbf{CUDA:} Versión 12.6.
    \item \textbf{Memoria RAM:} Gestión del flujo de puntos de colocación y
    ejecución del optimizador híbrido.
\end{itemize}

\subsection{Ecosistema de software}

\begin{itemize}
    \item \textbf{Python 3.11:} Lenguaje de implementación principal.
    \item \textbf{PyTorch 2.4:} Framework de \textit{Deep Learning} que
    proporciona diferenciación automática de segundo orden y soporte GPU
    mediante CUDA.
    \item \textbf{NumPy / SciPy:} Operaciones numéricas y muestreo LHS
    (\texttt{scipy.stats.qmc.LatinHypercube}).
    \item \textbf{Matplotlib:} Visualización de campos, distribuciones
    estadísticas y curvas de entrenamiento.
    \item \textbf{Módulos propios:} \texttt{src/models.py} (arquitectura SIREN),
    \texttt{src/losses.py} (función de pérdida PINN), \texttt{src/utils.py}
    (métricas y visualización).
\end{itemize}
